\section{协议噪声：Prompt、Seed、Decode 为什么会改变结论}

\subsection{Prompt template 不是中立变量}
LLM 任务里，prompt 经常携带隐含 inductive bias。即便任务语义相同，不同模板也可能导致明显分差。Sida 在讲中强调，如果论文结论依赖某一个模板，必须检查它在模板扰动下是否稳定。

\begin{knowledgebox}{协议敏感性测试（Protocol Sensitivity Test）}
固定模型与数据，仅改变模板/解码参数，观察分数变化范围。若变化过大，说明结论脆弱。
\end{knowledgebox}

\subsection{Decode 设置与 variance}
许多 benchmark 默认温度为 0，但在 agentic workload 中，模型通常无法完全 deterministic。此时 seed 和 decode policy 的影响会累积到轨迹层。课程里建议把 decode policy 当作评测协议的一部分进行 versioning，而不是隐含默认值。

\begin{importantbox}{协议版本管理建议}
\begin{itemize}
    \item 固定并公开：temperature、top-p、max tokens、stop criteria。
    \item 记录并公开：evaluation harness commit hash。
    \item 对关键结果做 seed sweep，至少报告 3-5 个随机种子的分布。
\end{itemize}
\end{importantbox}

\subsection{Judge 噪声：LLM-as-judge 的双刃剑}
LLM judge 能快速扩展评测规模，但会引入 judge drift、position bias、verbosity bias 等问题。Sida 的态度是 ``可用，但要校准''。常见做法包括双 judge 交叉验证、人评抽检、以及对评审 rubric 进行对抗测试。

\begin{warningbox}{Judge 过拟合风险}
当模型针对 judge 偏好优化（而非任务真实目标）时，会出现 ``会得分但不会做事'' 的现象。
\end{warningbox}

\subsection{本章小结}
协议层噪声是最容易被忽略、也最容易被治理的一层。只要建立版本化协议、重复评测和 judge 校准流程，很多 ``神秘波动'' 都能被解释为协议差异，而非能力突变。

